\documentclass[a4paper,11pt]{article}
\usepackage[utf8]{inputenc}
\usepackage[polish]{babel}
\usepackage{graphicx}
\usepackage{fancyhdr}
\usepackage{lastpage}
\usepackage{listings}
\usepackage{xcolor}
\usepackage{float}
\usepackage[margin=2cm]{geometry}

\lstset{
	basicstyle=\ttfamily\small,
	breaklines=true,
	frame=single,
	backgroundcolor=\color{gray!10},
	numbers=left,
	numberstyle=\tiny,
	keywordstyle=\color{blue}\bfseries,
	commentstyle=\color{green!60!black},
	stringstyle=\color{red}
}

\title{Laboratorium nr 5\\
	\large Usuwanie kolumny z ograniczeniami}
\author{Paweł Myszka, nr albumu 331720}
\date{13 grudnia 2025}

\pagestyle{fancy}
\fancyhf{}
\fancyhead[L]{Z5 -- 331720}
\fancyhead[R]{Administracja BD}
\fancyfoot[C]{\thepage}

\begin{document}
	
	\maketitle
	
	\section{Zadanie}
	
	Napisać procedurę \texttt{usun\_kol}, która automatycznie usunie kolumnę z tabeli wraz ze \textbf{wszystkimi} jej ograniczeniami (CONSTRAINTS). Procedura powinna:
	\begin{itemize}
		\item Przyjmować parametry: nazwa\_bazy, nazwa\_tabeli, nazwa\_kolumny
		\item Znaleźć wszystkie constraints związane z daną kolumną (DEFAULT, FOREIGN KEY, PRIMARY KEY, itp.)
		\item Usunąć constraints w pętli/kursorze
		\item Usunąć kolumnę
		\item Wyświetlić dowód działania (SELECT przed i po usunięciu)
	\end{itemize}

	\section{Problem}
	
	Nie zawsze jest możliwe bezpośrednie usunięcie kolumny, szczególnie gdy kolumna posiada przypisane ograniczenia. Próba wykonania:
	
	\begin{lstlisting}[language=SQL]
ALTER TABLE DETAIL_TB DROP COLUMN MASTER_ID
	\end{lstlisting}
	
	powoduje błąd:
	\begin{lstlisting}
Msg 5074: The object 'DF__DETAIL_TB__maste__3E52440B' is dependent on column 'MASTER_ID'.
Msg 5074: The object 'FK_DETAIL__MASTER_TB' is dependent on column 'MASTER_ID'.
Msg 4922: ALTER TABLE DROP COLUMN MASTER_ID failed because one or more objects access this column.
	\end{lstlisting}
	
	\section{Implementacja}
	
	\subsection{Procedura usun\_kol}
	
	\begin{lstlisting}[language=SQL]
CREATE PROCEDURE usun_kol
	@nazwa_bazy NVARCHAR(128),
	@nazwa_tabeli NVARCHAR(128),
	@nazwa_kolumny NVARCHAR(128)
AS
BEGIN
	DECLARE @sql NVARCHAR(MAX)
	DECLARE @constraint_name NVARCHAR(128)
	
	-- Sprawdzenie czy kolumna istnieje
	IF NOT EXISTS (
		SELECT 1 FROM INFORMATION_SCHEMA.COLUMNS
		WHERE TABLE_CATALOG = @nazwa_bazy
		AND TABLE_NAME = @nazwa_tabeli
		AND COLUMN_NAME = @nazwa_kolumny
	)
	BEGIN
		RAISERROR('Kolumna %s nie istnieje w tabeli %s.%s', 
			16, 1, @nazwa_kolumny, @nazwa_bazy, @nazwa_tabeli)
		RETURN
	END
	
	-- Kursor do usunięcia wszystkich constraints związanych z kolumną
	DECLARE constraint_cursor CURSOR FOR
		SELECT CONSTRAINT_NAME
		FROM INFORMATION_SCHEMA.CONSTRAINT_COLUMN_USAGE
		WHERE TABLE_CATALOG = @nazwa_bazy
		AND TABLE_NAME = @nazwa_tabeli
		AND COLUMN_NAME = @nazwa_kolumny
		UNION
		SELECT d.name
		FROM sys.default_constraints d
		JOIN sys.columns c ON d.parent_object_id = OBJECT_ID(@nazwa_tabeli)
		AND d.parent_column_id = c.column_id
		WHERE c.name = @nazwa_kolumny
	
	OPEN constraint_cursor
	FETCH NEXT FROM constraint_cursor INTO @constraint_name
	
	WHILE @@FETCH_STATUS = 0
	BEGIN
		SET @sql = N'ALTER TABLE [' + @nazwa_bazy + '].dbo.[' 
			+ @nazwa_tabeli + '] DROP CONSTRAINT [' 
			+ @constraint_name + ']'
		EXEC sp_executesql @sql
		PRINT 'Usunięto constraint: ' + @constraint_name
		
		FETCH NEXT FROM constraint_cursor INTO @constraint_name
	END
	
	CLOSE constraint_cursor
	DEALLOCATE constraint_cursor
	
	-- Usunięcie kolumny
	SET @sql = N'ALTER TABLE [' + @nazwa_bazy + '].dbo.[' 
		+ @nazwa_tabeli + '] DROP COLUMN [' + @nazwa_kolumny + ']'
	EXEC sp_executesql @sql
	PRINT 'Usunięto kolumnę: ' + @nazwa_kolumny
	
END
	\end{lstlisting}
	
	\section{Wyjaśnienie Implementacji}
	
	\subsection{Sprawdzenie kolumny}
	
	Procedura najpierw weryfikuje, czy dana kolumna istnieje w tabeli, korzystając z widoku \texttt{INFORMATION\_SCHEMA.COLUMNS}. Jeśli kolumna nie istnieje, procedura zwraca błąd za pomocą \texttt{RAISERROR}.
	
	\subsection{Znajdowanie Constraints}
	
	Aby znaleźć wszystkie constraints powiązane z kolumną, procedura korzysta z dwóch źródeł:
	
	\begin{itemize}
		\item \texttt{INFORMATION\_SCHEMA.CONSTRAINT\_COLUMN\_USAGE} -- zawiera FOREIGN KEY, PRIMARY KEY, UNIQUE constraints
		\item \texttt{sys.default\_constraints} -- zawiera DEFAULT constraints
	\end{itemize}
	
	Obie tabele są połączone za pomocą \texttt{UNION}, aby uzyskać kompletną listę.
	
	\subsection{Usuwanie Constraints}
	
	Kursor iteruje po wszystkich znalezionych constraints i usuwa je jeden po drugim przy użyciu:
	
	\begin{lstlisting}[language=SQL]
ALTER TABLE [nazwa_bazy].dbo.[nazwa_tabeli] DROP CONSTRAINT [nazwa_constraint]
	\end{lstlisting}
	
	Dla każdego usuniętego constraintu wyświetlony jest komunikat \texttt{PRINT}.
	
	\subsection{Usuwanie Kolumny}
	
	Po usunięciu wszystkich constraints, procedura usuwa samą kolumnę:
	
	\begin{lstlisting}[language=SQL]
ALTER TABLE [nazwa_bazy].dbo.[nazwa_tabeli] DROP COLUMN [nazwa_kolumny]
	\end{lstlisting}
	
	\section{Testowanie}
	
	\subsection{Przygotowanie}
	
	Procedura została przetestowana na dwóch tabelach:
	\begin{itemize}
		\item \texttt{MASTER\_TB} -- zawiera klucz główny
		\item \texttt{DETAIL\_TB} -- zawiera kolumnę \texttt{MASTER\_ID} z dwoma constraints:
		\begin{itemize}
			\item DEFAULT constraint (wartość domyślna: M01)
			\item FOREIGN KEY constraint (odniesienie do MASTER\_TB)
		\end{itemize}
	\end{itemize}
	
	\subsection{Wynik Testowania}
	
	Procedura \texttt{usun\_kol} została pomyślnie uruchomiona:
	
	\begin{lstlisting}
EXEC usun_kol 
	@nazwa_bazy = 'B25_ADM_PM331720', 
	@nazwa_tabeli = 'DETAIL_TB', 
	@nazwa_kolumny = 'MASTER_ID'
	\end{lstlisting}
	
	Wynik:
	\begin{lstlisting}
Usunięto constraint: DF__DETAIL_TB__maste__3D2915A8
Usunięto constraint: FK_DETAIL__MASTER_TB
Usunięto kolumnę: MASTER_ID
	\end{lstlisting}
	
	Przed usunięciem kolumna \texttt{MASTER\_ID} zawierała wartość:
	\begin{lstlisting}
master_id            DETAIL_ID
-------------------- --------------------
M01                  D01
	\end{lstlisting}
	
	Po usunięciu kolumny tabela zawiera tylko:
	\begin{lstlisting}
DETAIL_ID
--------------------
D01
	\end{lstlisting}
	
	\section{Wnioski}
	
	\begin{itemize}
		\item ✓ Procedura \texttt{usun\_kol} prawidłowo identyfikuje wszystkie constraints powiązane z kolumną
		\item ✓ Constraints (zarówno DEFAULT jak i FOREIGN KEY) są poprawnie usuwane
		\item ✓ Kolumna jest usuwana dopiero po usunięciu wszystkich constraints
		\item ✓ Procedura wyświetla komunikaty potwierdzające każdy krok operacji
		\item ✓ Procedura zabezpiecza się przed usunięciem nieistniejącej kolumny
		\item ✓ Rozwiązanie umożliwia automatyzację procesu usuwania kolumn z ograniczeniami
	\end{itemize}
	
\end{document}
